%% kapa: exam=1 round=37 year=2026 period=2 form=A subject=appraisal_law subject_ko=감정평가관계법규 seq=26 number=26
%% source: Q-Net 감정평가사 기출문제 (cst00309) artlSeq=5258924 file=raw/1008/5258924/2026년 제37회 감정평가사 1차 2교시.pdf
%% conversion: kordoc 4.24.51 -> tools/gen1.py (auto); review=proofread-sonnet; flags=none
\begin{question}{26}
건축법령상 면적 등의 산정방법에 관한 조문의 일부이다. ( )에 들어갈 내용으로 옳은 것은?
\begin{framedtext}
3. ( ㄱ ): 건축물의 각 층 또는 그 일부로서 벽, 기둥, 그 밖에 이와 비슷한 구획의 중심선으로 둘러싸인 부분의 수평투영면적으로 한다. 다만, 다음 각 목의 어느 하나에 해당하는 경우에는 각 목에서 정하는 바에 따른다. \par
\hspace*{1em}가. 벽ㆍ기둥의 구획이 없는 건축물은 그 지붕 끝부분으로부터 수평거리 ( ㄴ )미터를 후퇴한 선으로 둘러싸인 수평투영면적으로 한다. \par
\end{framedtext}
\begin{choices}
  \item ㄱ: 바닥면적, ㄴ: 0.5
  \item ㄱ: 바닥면적, ㄴ: 1
  \item ㄱ: 바닥면적, ㄴ: 1.5
  \item ㄱ: 건축면적, ㄴ: 1
  \item ㄱ: 건축면적, ㄴ: 1.5
\end{choices}
\end{question}
\answer{2}
